% !TEX root = tarasov-qasp-zh.tex

\chapter[你能确定摩擦力吗？]{你能确定摩擦力吗？\label{ch-03}}

\begin{dialogue}
\tchr 我想更详细地谈谈各种问题中摩擦力的计算。我指的是干滑动摩擦（dry sliding friction）——两滑动表面之间若没有任何物质层（例如润滑剂），这样的滑动摩擦就称为干摩擦。

\std 可是这里的一切看起来都相当清楚呀。
\begin{figure}[!ht]
\centering
\includegraphics[width=0.5\textwidth]{figs/ch-03/fig-015a.pdf}
\caption[用绳拉动雪橇。]{用绳拉动雪橇。\label{fig-15}}
\end{figure}
\tchr 尽管如此，考试中出的许多错误恰恰源于不会计算摩擦力。来看\figr{fig-15} 所示的例子。\hlred{重力为 $P$ 的雪橇被力 $F$ 拉动，力施加在与水平方向成 $\alpha$ 角的绳上；摩擦系数为 $k$。求滑动摩擦力。}你打算怎样入手？

\std 哎呀，这个题目看起来非常简单。摩擦力等于 $kP$。

\tchr 完全不对。滑动摩擦力并不等于 $kP$，而等于 $kN$，其中 $N$ 是支承力。回忆一下\chap{ch-02} 中的式~\eqref{eq-05}。

\std 可是，这难道不是一回事吗？

\tchr 在个别情形下，重力与支承力可能彼此相等，但一般说来，它们是截然不同的两个力。就拿我刚才提出的例子来看。作用在物体（雪橇）上的力是：重力 $P$、支承力 $N$、滑动摩擦力 $F_{f\!\!r}$ 和绳的张力 $F$（见\figr{fig-15}）。把力 $F$ 分解为竖直分量（$F \sin \alpha$）与水平分量（$F \cos \alpha$）。所有竖直方向上的力相互平衡，由此便可以求出支承力：
\begin{equation}%
N = P - F \sin \alpha
\label{eq-06}
\end{equation}
如你所见，这个力不等于雪橇的重力，而是小了 $F \sin \alpha$。从物理上看，这是意料之中的：绷紧的绳斜向上拉，多少把雪橇往上“提”了一些。这就减小了雪橇压在表面上的力，支承力也就随之减小。因此，本题中
\begin{equation}%
F_{f\!\!r}=k(P - F \sin \alpha)
\label{eq-07}
\end{equation}
如果绳是水平的（$\alpha=0$），那么代替式~(\ref{eq-06}) 的将是 $N=P$，由此可得 $F_{f\!\!r}=kP$。
\pnote{错误根源：把摩擦定律中的正压力当成了重力。人教版必修第一册第三章“摩擦力”一节写作 $F=\mu F_{N}$，正压力 $F_{N}$（即原书的 $N$）一般不等于重力，须由垂直于接触面方向的平衡条件求出：斜向上拉得 $N=P-F\sin\alpha$；若改为斜向下推，则 $N=P+F\sin\alpha$，摩擦力反而更大。}

\std 现在我明白了。这一点我以前从来没有想过。

\tchr 这是考生中相当普遍的一种错误：把滑动摩擦力当作摩擦系数与重力的乘积，而不是与支承力的乘积。今后要力求避免这类错误。

\std 我要遵循这样一条规则：\hlblue{求摩擦力，先求支承力。}

\tchr 到现在为止，我们讨论的都是滑动摩擦力。现在来看静摩擦力（static friction）。它有一些学生未必充分注意的特点。看下面这个例子。物体静止在水平面上，受到一个水平力 $F$ 的作用，此力有使物体移动的趋势。你认为在这种情形下摩擦力有多大？

\std 如果物体静止在水平面上，而力 $F$ 沿水平方向作用，那么 $N=P$。对吗？

\tchr 完全正确。继续。

\std 由此可得，摩擦力等于 $kP$。
\begin{figure}[!ht]%
\centering
\includegraphics[width=0.45\textwidth]{figs/ch-03/fig-016a.pdf}
\caption[作用在静止物体上的力。]{作用在静止物体上的力。\label{fig-16}}
\end{figure}
\tchr 你犯了一个典型错误：混淆了滑动摩擦力与静摩擦力。假如物体是在沿平面滑动，你的答案就对了。但这里物体是静止的，因此，作用在物体上的所有力必须相互平衡。物体受到四个力：重力 $P$、支承力 $N$、力 $F$ 和静摩擦力 $F_{f\!\!r}$（\figr{fig-16}）。竖直方向的力 $P$ 与 $N$ 相互平衡，水平方向的力 $F$ 与 $F_{f\!\!r}$ 也应当相互平衡。于是
\begin{equation}%
F_{f\!\!r} = F
\label{eq-08}
\end{equation}
\pnote{初学者最常踩的坑正在这里：物体未滑动时，静摩擦力绝不能写成 $kN$，只能由平衡条件求出（此处 $F_{f\!\!r}=F$），取值范围为 $0 \le F_{f\!\!r} \le k_{0}N$。人教版必修第一册第三章“摩擦力”一节的图像与此相同：静摩擦力随外力增大而增大，至最大静摩擦力为止，物体开始滑动后才转为滑动摩擦力。}
\std 由此可见，静摩擦力取决于那个试图移动物体的外力。

\tchr 是的，正是如此。静摩擦力随力 $F$ 的增大而增大，但它不会无限增大。静摩擦力会达到一个最大值：
\begin{equation}%
F_{max} = k_{0}N
\label{eq-09}
\end{equation}


系数 $k_{0}$ 稍大于系数 $k$——按式~(\ref{eq-05})，$k$ 表征的是滑动摩擦力。一旦外力 $F$ 达到值 $k_{0}N$，物体就开始滑动。此时系数 $k_{0}$ 变得与 $k$ 相等，因而摩擦力略有减小。力 $F$ 再继续增大时，摩擦力（此时已是滑动摩擦力）便不再进一步增大（直到达到很高的速度之前），物体以逐渐增大的加速度运动。许多考生不会确定摩擦力这一点，用下面一个相当简单的问题就能暴露出来：\hlred{重力为 $P$ 的物体静止在倾角为 $\alpha$ 的斜面上，摩擦力是多大？}人们会听到各种各样的错误答案：有人说摩擦力等于 $kP$，也有人说它等于 $kN =kP \cos \alpha$。

\std 我明白了。既然物体静止，我们要处理的就是静摩擦力。它应当由沿斜面方向作用的各力的平衡条件来求。在本题中，这样的力有两个：摩擦力 $F_{f\!\!r}$ 和沿斜面向下作用的下滑力 $P \sin \alpha$。因此，正确的答案是
\begin{equation*}%
F_{f\!\!r}=P \sin \alpha
\end{equation*}
\pnote{注意答案与摩擦系数无关：只要物体静止，静摩擦力就由平衡条件确定。物体能静止于斜面的条件是 $P\sin\alpha \le k_{0}P\cos\alpha$，即 $\tan\alpha \le k_{0}$（$k_{0}$ 为最大静摩擦力与正压力之比，即静摩擦系数）。临界倾角 $\varphi=\arctan k_{0}$ 称为摩擦角（angle of friction），中学不作要求，物理竞赛中常用于判断斜面上的物体能否自锁。}
\tchr 完全正确。最后，考虑\figr{fig-17} 所示的问题。\hlred{质量为 $m$ 的重物放在质量为 $M$ 的物体上；两者之间的最大静摩擦力由系数 $k_{0}$ 表征，物体与地面之间没有摩擦。求施于物体的最小力 $F$，在这个力作用下重物将开始沿物体滑动。}
\begin{figure}[!h]%
\centering
\includegraphics[width=0.4\textwidth]{figs/ch-03/fig-017a.pdf}
\caption[重物叠放在物体上，求使其开始滑动的最小力。]{质量为 $m$ 的重物放在质量为 $M$ 的物体上。须求使重物开始沿物体滑动的最小力 $F$。\label{fig-17}}
\end{figure}
\std 我先假定力 $F$ 足够小，重物不会沿物体滑动。这时两个物体将获得加速度
\begin{equation*}%
a= \frac{F}{M+m}
\end{equation*}

\tchr 正确。那么，是什么力使重物获得这个加速度的呢？

\std 重物将受到静摩擦力 $F_{f\!\!r}$ 的作用，正是它使重物产生这个加速度。于是
\begin{equation*}%
F_{f\!\!r} = ma= \frac{Fm}{M+m}
\end{equation*}
由此可见，随着力 $F$ 的增大，静摩擦力 $F_{f\!\!r}$ 也增大。但它不能无限增大，其最大值为
\begin{equation*}%
F_{\textrm{fr \hairsp max}}= k_{0}N = k_{0}mg
\end{equation*}
因而，两个物体仍能作为一个整体一起运动时，力 $F$ 的最大值由条件
\begin{equation*}%
 k_{0}mg = \frac{Fm}{M+m}
\end{equation*}
确定。这也就是使重物开始沿物体滑动的最小力。
\pnote{底本章内另有几处排印小误（如“$F$ 与 $F_{f\!\!r}$”之后脱句点、一处冒号应为句点、“force $F$. at which”衍句点、\figr{fig-17} 图题``so that load begins to slide''脱``the''），均径改不注。本组对话底本英文语序错乱两处（``What force will this acceleration impart to the load?''与``It will be subjected to the force of static friction … by the acceleration.''），文意为“是什么力使重物获得这个加速度”，显系排印或俄译英之误，中文按文意通译。}
\pnote{使重物获得加速度的正是静摩擦力——摩擦力可以与运动方向相同、充当动力。由最后两式可解出所求的最小力 $F=k_{0}g(M+m)$，它与重物质量 $m$ 无关：$m$ 增大时，带动重物所需的摩擦力与最大静摩擦力按同一比例增大。这类叠放体问题即高考常见的“板块模型”，须整体法（求共同加速度）与隔离法（求静摩擦力）配合使用（参看\chap{ch-02} 评注），与人教版必修第一册第四章“牛顿运动定律的应用”一节的题型相当。}

\tchr 你对所提问题的解答是正确的。我对你的推理完全满意。

\end{dialogue}

% NEWTERMS: dry (sliding) friction | 干（滑动）摩擦 | 两滑动面间无润滑层等任何物质层
% NEWTERMS: lubricant | 润滑剂 | 普通词，供一致性备查
% NEWTERMS: static friction | 静摩擦（力） | force of static friction 译"静摩擦力"；人教版必修第一册第三章同名
% NEWTERMS: maximum force of static friction | 最大静摩擦力 | $F_{max}=k_0N$；人教版同名
% NEWTERMS: coefficient of friction | 摩擦系数 | 泛称；$k$ 为滑动摩擦系数，$k_0$ 为静摩擦系数（原书未给 $k_0$ 命名）
% NEWTERMS: sliding force | 下滑力 | 重力沿斜面方向的分力，沿第 2 章既有译名
% NEWTERMS: sled | 雪橇 | 普通词
% NEWTERMS: load | 重物 | 语境为叠放在物体上的负载
% NEWTERMS: angle of friction | 摩擦角 | 评注层引入；中学不要求，竞赛常用
